\documentclass{article}

\usepackage[margin=1in]{geometry}
\usepackage[T1]{fontenc}
\usepackage{mathptmx}
\usepackage{courier}
\usepackage{microtype}
\usepackage{algorithm}
\usepackage{algpseudocode}
\usepackage{amsmath}
\usepackage{xcolor}
\usepackage{listings}
\usepackage{float}

\definecolor{PromptBg}{RGB}{248,249,250}
\definecolor{PromptRule}{RGB}{210,214,220}

\lstdefinestyle{prompt}{
  basicstyle=\ttfamily\footnotesize,
  breaklines=true,
  breakatwhitespace=true,
  columns=fullflexible,
  frame=single,
  framerule=0.35pt,
  rulecolor=\color{PromptRule},
  backgroundcolor=\color{PromptBg},
  xleftmargin=0.6em,
  xrightmargin=0.6em,
  aboveskip=0.6em,
  belowskip=0.9em,
  keepspaces=true,
  showstringspaces=false
}

\lstset{style=prompt}

\begin{document}

\pagestyle{empty}

\begin{center}
{\Large\bfseries Filesystem Retrieval Flow}
\end{center}

\section*{1. Strategy Selection}

The filesystem retriever treats the repository as a tree. Each node is either a
directory or a file, and each candidate shown to the LLM contains only the node
id, relative path, and type. Metadata fields are reserved for a later version
and are not rendered in the active LLM-facing candidate payload.

The query entry point chooses the retrieval strategy in this order:
\begin{enumerate}
  \item Read the tree metadata. If the tree source is not filesystem, dispatch
  to the document retriever.
  \item If the caller passes \texttt{strategy="beam"}, use Beam directly.
  If the caller passes \texttt{strategy="block"}, use Block directly.
  \item If the caller passes \texttt{strategy="auto"}, count tree nodes.
  Trees with at most 50 nodes use Beam; larger trees use Block.
  \item Run the chosen retriever with the query, beam width, and final result
  limit.
\end{enumerate}

The algorithms below show the filesystem path. Candidate order follows the
stored filesystem order except for the optional post-merge path ordering in
Algorithm 5.

\section*{2. Beam Retriever}

Beam retrieval is level-by-level tree navigation. On each turn, every current
beam is expanded to its direct children. The union of those children is the
candidate set for that LLM call. If a beam is already a leaf, the leaf itself is
kept as a candidate so it can be returned.

The Beam flow is implemented by Algorithm 1:
\begin{enumerate}
  \item Initialize the beam set with the root node.
  \item Expand every current beam into direct children.
  \item Treat the union of expanded children as the current candidate set.
  \item Ask the LLM to rank candidates and optionally signal \texttt{done}.
  \item If the LLM says \texttt{done} but returned only directories, force
  \texttt{done=false} and continue exploring.
  \item Add returned ids to the running top-candidate set.
  \item Keep the top ranked expandable ids as the next beams.
  \item Stop when $k$ file results have been collected, the LLM says
  \texttt{done}, there are no candidates, or the depth limit is reached.
\end{enumerate}

\begin{algorithm}[H]
\caption{FilesystemBeam}
\begin{algorithmic}[1]
\Require Tree id $T$, query $q$, beam width $b$, final result limit $k$
\Ensure Retrieved file nodes $R$
\State $r \gets \Call{Root}{T}$
\State $B \gets [r]$
\State $A \gets [\ ]$
\State $p \gets \max(b,k)$
\For{$t = 0,1,\ldots,\Call{MaxDepth}{T}$}
  \State $C \gets [\ ]$
  \ForAll{$u \in B$}
    \State $Ch \gets \Call{Children}{T,u}$
    \If{$Ch = \emptyset$}
      \State $C \gets C \cup \{u\}$
    \Else
      \State $C \gets C \cup Ch$
    \EndIf
  \EndFor
  \If{$C = \emptyset$}
    \State \textbf{break}
  \EndIf
  \State $(L, done) \gets \Call{LLMRankBeam}{q, C, A, p}$
  \If{$done$ and no node in $L$ is a file}
    \State $done \gets \textsc{false}$
  \EndIf
  \State $Q \gets [u \in L: u \text{ is a file}]$
  \State $A \gets (A \cup Q)[1:k]$
  \If{$|A| \ge k$}
    \State \textbf{break}
  \EndIf
  \State $B \gets [\ ]$
  \ForAll{$u \in L$}
    \If{$u$ is a directory or expandable node}
      \State append $u$ to $B$
    \EndIf
    \If{$|B| = b$}
      \State \textbf{break}
    \EndIf
  \EndFor
  \If{$done$}
    \State \textbf{break}
  \EndIf
\EndFor
\State \Return $A[1:k]$
\end{algorithmic}
\end{algorithm}

\section*{3. Block Retriever}

Block retrieval has the same recursive intent as Beam, but it does not present
one large candidate list when the directory is wide. Instead, it partitions the
current visible children into token-bounded blocks. For a block $B_i$, the block
candidate set is exactly $C_i = B_i.\text{\texttt{node\_ids}}$. The LLM returns a
block-local ordered list from that block only; those returned ids are not
inserted back into $C_i$. They are merged with other block results by order
preservation, then split into file results and directory frontier nodes.
When all entries in a block share a directory prefix, the rendered block may add
\texttt{Path prefix: <prefix>} once and show each candidate path relative to that
prefix.

The Block flow is implemented by Algorithms 2--5:
\begin{enumerate}
  \item Initialize the frontier with the root, then collapse virtual
  single-child directory chains. This is handled by the main loop in
  Algorithm 2.
  \item Materialize visible candidates from the current frontier. At turn 0 this
  is the top visible layer; later turns use the direct children of frontier
  nodes. This is handled by Algorithm 3.
  \item Pack visible candidates into token-bounded blocks
  $\mathcal{B}=\{B_1,\ldots,B_m\}$. This is handled by Algorithm 3.
  \item Render each block using only id, path, and type. If a shared path prefix
  reduces tokens, render the prefix once and shorten candidate paths within that
  block. This is handled by Algorithm 3.
  \item For each block $B_i$, define the allowed candidate set
  $C_i=B_i.\text{\texttt{node\_ids}}$ and ask the LLM to rank only ids in $C_i$.
  This is handled by Algorithm 4.
  \item Keep per-block outputs as block-local groups
  $\mathcal{G}=[L_1,\ldots,L_m]$ and derive
  $M=\Call{Flatten}{\mathcal{G}}$ for the no-ranker path. This merge removes
  duplicates while preserving each block's LLM-returned order.
  \item Optionally reorder across block groups without changing the relative
  order inside any $L_i$, then split the resulting list by node type: files enter
  the accumulated top-candidate result set, while directories enter the next
  frontier. This is handled by Algorithm 5.
  \item If the LLM says \texttt{done} but the frontier still contains
  directories, override \texttt{done=false} and continue drilling down. This is
  handled by Algorithm 5.
  \item Stop immediately once the accumulated file result set reaches the final
  result limit $k$.
\end{enumerate}

\begin{algorithm}[H]
\caption{FilesystemBlock Main Loop}
\begin{algorithmic}[1]
\Require Tree id $T$, query $q$, beam width $b$, final result limit $k$, token budget $\tau$
\Ensure Retrieved file nodes $R$
\State $r \gets \Call{Root}{T}$
\State $F \gets \Call{CollapseSingleChildDirs}{[r]}$
\State $A \gets [\ ]$
\State $P \gets [\ ]$
\State $\hat b \gets b$ if set, otherwise $1$
\State $p \gets \max(\hat b,k)$ \Comment{LLM pick limit, not the final result limit}
\State $t \gets 0$
\While{$t < \Call{MaxRounds}{}$}
  \If{$t > 0$ and no node in $F$ has children}
    \State \textbf{break}
  \EndIf
  \State $(\mathcal{B}, X) \gets \Call{BuildBlockCandidates}{T,F,t,\tau}$
  \If{$X = \emptyset$ or $\mathcal{B}=\emptyset$}
    \State \textbf{break}
  \EndIf
  \State $p_t \gets p$ if $|\mathcal{B}|=1$, otherwise $\max(p,2)$
  \State $(\mathcal{G}, M, done) \gets \Call{CallBlocksAndMerge}{q,\mathcal{B},F,P,p_t}$
  \State $(F,P,A,done) \gets \Call{UpdateState}{T,q,\mathcal{G},M,done,b,k,A}$
  \If{$|A| \ge k$ or $done$}
    \State \textbf{break}
  \EndIf
  \State $t \gets t + 1$
\EndWhile
\State \Return $A[1:k]$
\end{algorithmic}
\end{algorithm}

\begin{algorithm}[H]
\caption{BuildBlockCandidates}
\begin{algorithmic}[1]
\Require Tree id $T$, frontier $F$, turn $t$, token budget $\tau$
\Ensure Token-bounded blocks $\mathcal{B}$ and visible candidates $X$
\If{$t = 0$}
  \State $X \gets \Call{TopVisibleNodes}{T,F_1}$
\Else
  \State $X \gets [\ ]$
  \ForAll{$u \in F$}
    \State $X \gets X \cup \Call{Children}{T,u}$
  \EndFor
\EndIf
\If{$X = \emptyset$}
  \State \Return $(\emptyset,\emptyset)$
\EndIf
\State Render each candidate using id, path, and type
\State Render a shared \texttt{Path prefix} when it shortens the block
\State $\mathcal{B} \gets \Call{PackIntoBlocks}{X,\tau}$
\State \Return $(\mathcal{B},X)$
\end{algorithmic}
\end{algorithm}

\begin{algorithm}[H]
\caption{CallBlocksAndMerge}
\begin{algorithmic}[1]
\Require Query $q$, blocks $\mathcal{B}$, frontier $F$, previous top candidates $P$, LLM pick limit $p$
\Ensure Block-local id groups $\mathcal{G}$, merged ids $M$, and stop signal $done$
\ForAll{$i=1,\ldots,|\mathcal{B}|$ \textbf{in parallel}}
  \State $B_i \gets \mathcal{B}_i$
  \State $C_i \gets B_i.\text{\texttt{node\_ids}}$
  \State $(L_i,d_i) \gets \Call{LLMRankBlock}{q,B_i,C_i,F,P,p}$
\EndFor
\State $\mathcal{G} \gets [\ ]$
\For{$i=1,\ldots,|\mathcal{B}|$}
  \State $\mathcal{G} \gets \mathcal{G} \parallel [L_i]$
\EndFor
\State $\mathcal{G} \gets \Call{UniquePreserveBlockOrder}{\mathcal{G}}$
\State $M \gets \Call{Flatten}{\mathcal{G}}$
\State $done \gets \Call{AllTrue}{\{d_i\}}$
\State \Return $(\mathcal{G},M,done)$
\end{algorithmic}
\end{algorithm}

\begin{algorithm}[H]
\caption{UpdateState}
\begin{algorithmic}[1]
\Require Tree id $T$, query $q$, block-local id groups $\mathcal{G}$, merged ids $M$, stop signal $done$, beam width $b$, final result limit $k$, top candidates $A$
\Ensure Next frontier $F$, previous top candidates $P$, top candidates $A$, stop signal $done$
\State $M' \gets \Call{ConstrainedPathOrder}{\mathcal{G},q}$
\State $Q \gets [u \in M': u \text{ is a file}]$
\State $D \gets [u \in M': u \text{ is a directory}]$
\State $r \gets \max(0, k-|A|)$
\State $A \gets A \cup Q[1:r]$
\State $P \gets A[1:k]$
\If{$|A| \ge k$}
  \State \Return $(\emptyset,P,A[1:k],\textsc{true})$
\EndIf
\State $F' \gets D[1:b]$
\State $F \gets \Call{CollapseSingleChildDirs}{F'}$
\If{$done$ and some node in $F$ still has children}
  \State $done \gets \textsc{false}$
\EndIf
\State \Return $(F,P,A[1:k],done)$
\end{algorithmic}
\end{algorithm}

\paragraph{Block semantics.}
For every block call, the LLM can only select ids from the current block's
allowed list. If the current directory produces three blocks
$B_1,B_2,B_3$, the calls produce $L_1,L_2,L_3$. The algorithm keeps these as
block-local groups $\mathcal{G}=[L_1,L_2,L_3]$ and derives
$M=\Call{Flatten}{\mathcal{G}}$. This is a merge, not another LLM ranking step.
$M$ is therefore called \texttt{merged\_ids}, not a ranked list. Without path
ordering, its order is inherited from block order and each block's local LLM
order. With optional path ordering, only cross-block order may change; the
relative order inside each $L_i$ is fixed. Files enter
\texttt{top\_candidate\_ids}; directories enter the next frontier after
beam-width truncation. The next candidate set is produced only by expanding that
frontier.

\paragraph{Block cache method.}
Block retrieval separates stable tree text from dynamic query state. The stable
part is the rendered block content; the dynamic part contains the query, current
exploration path, previous top candidates, and the ids allowed for the current
block call. If cache is enabled, the retriever builds a cache payload from a
static instruction segment plus a bounded cache window of previously useful
block segments. The current block is also placed in the cache payload when
\texttt{cache\_current\_block} is enabled at the top level, or when
\texttt{cache\_subtree\_block} is enabled in recursive subtree rounds.

After each round, the blocks that contain the next frontier are appended to the
cache window. The first top-level block may be pinned, while later entries can be
trimmed when the cache-window token budget is exceeded. The LLM call then sends
the cache payload as cached content and sends only the query-specific message as
the dynamic user message. If the LLM adapter does not support cached calls, the
same content is concatenated into a normal prompt; retrieval semantics are the
same, but no provider-side cache saving is used. The \texttt{Path prefix} line
is separate from this cache mechanism: it is only a token-compression rendering
trick inside a block.

\paragraph{Path ordering.}
The deterministic merge preserves block order and block-local order. It does not
compare candidates from different blocks. When a ranker is supplied, the path
scorer runs once over the block-local groups
$\mathcal{G}=[L_1,\ldots,L_m]$ and before the file/directory split:
\[
M'=\Call{ConstrainedPathOrder}{\mathcal{G},q}.
\]
Only the current head of each group may be selected at each step, so different
blocks can be interleaved while the relative order inside every $L_i$ remains
unchanged. Without a ranker, $M'=\Call{Flatten}{\mathcal{G}}$.

The optional scorer is path-only:
\[
s(u,q)=3\,\mathrm{BM25}(q,\mathrm{basename}(u))
      +1.5\,\mathrm{BM25}(q,\mathrm{parent}(u))
      +\mathrm{BM25}(q,\mathrm{path}(u)).
\]
The ordered list $M'$ is then split: files fill the remaining $k-|A|$ result
slots, and directories are truncated to the beam width for the next frontier.

\section*{4. Output Semantics}

In filesystem mode, \texttt{candidate\_set}, \texttt{ranked\_ids},
\texttt{frontier}, and \texttt{top\_candidate\_ids} have different meanings.
The candidate set is the current block-local pool shown to the LLM.
\texttt{ranked\_ids} is only the LLM-returned order inside one block and may
contain files or directories. After deterministic merge, the cross-block list is
\texttt{merged\_ids}; it is not a global ranking. If optional path ordering is
enabled, the retriever derives one ordered list from \texttt{merged\_ids} before
splitting it. Files then enter \texttt{top\_candidate\_ids}, which is the
maintained result set, and directories form the next frontier. Final filesystem
output returns at most $k$ entries from \texttt{top\_candidate\_ids} and their
contents. The separate \texttt{pick\_limit} controls how many ids an LLM call may
return inside one round; it is not the final result count.

\newpage

\section*{Appendix A. LLM Prompt Payloads}

\subsection*{A.1 Tool schema}

Both retrievers use the same tool-call contract. The LLM must return one
\texttt{rank} tool call.

\noindent\textbf{Tool schema}
\begin{lstlisting}
TOOLS = [
    {
        "name": "rank",
        "description": "Rank candidate node ids for the query",
        "input_schema": {
            "type": "object",
            "properties": {
                "ranked_ids": {
                    "type": "array",
                    "items": {"type": "string"}
                },
                "done": {"type": "boolean"},
            },
            "required": ["ranked_ids"],
        },
    }
]
\end{lstlisting}

\subsection*{A.2 Rendered prompt examples}

The boxes below show examples after the templates are rendered for an LLM call.
Angle-bracket fields are filled from the current tree state before each call.

\subsection*{A.2.1 Beam prompt}

\noindent Source template: \texttt{contextdb/prompts/beam\_fs.jinja}
\begin{lstlisting}
You are navigating a source code repository to find files relevant to this query:

Query: <query>

Already top candidates:
- <previous top candidate node id>

Current candidates (pick up to <pick_limit>, most relevant first):
- id: <node id>
  path: <relative path>
  type: <directory | file>
- id: <node id>
  path: <relative path>
  type: <directory | file>

RULES:
- Select directories to explore deeper, or files if they match the query
- set done=true ONLY when you've found specific source files that answer the query
- set done=false if you returned directories that need further exploration

Return ONE tool call "rank" with:
- ranked_ids: list of candidate ids in best-to-worst order
- done: true or false
\end{lstlisting}

\subsection*{A.2.2 Block prompt}

\noindent Source templates: \texttt{block\_fs\_cache\_prefix.jinja} and
\texttt{block\_fs.jinja}. The first part is the cached block prefix; the second
part is the dynamic user message for the current query and allowed ids.

\noindent\textbf{Cached block prefix}
\begin{lstlisting}
=== REPOSITORY DIRECTORY TREE ===

Path prefix: <shared directory prefix>/

- id: n1
  path: <path relative to prefix>
  type: <directory | file>
- id: n2
  path: <path relative to prefix>
  type: <directory | file>

You are navigating a source code repository's directory tree to find files relevant to a query.
Each entry above has:
- id: a unique node identifier
- path: relative file/directory path
- type: file or directory

RULES:
- Select directories to drill deeper into, or files if they directly match the query
- set done=false if returned directories need further exploration
- Select ONLY ids from the allowed list provided below

Return ONE tool call "rank" with:
- ranked_ids: list of candidate ids in best-to-worst order
- done: true ONLY when the returned ids are specific source files
\end{lstlisting}

\noindent\textbf{Dynamic message}
\begin{lstlisting}
Query: <query>

Top candidates so far:
- <previous top candidate path>

Current exploration path:
- <current frontier path>

The LAST BLOCK above is the CURRENT BLOCK.
IDs in each block are local to that block.
Return only ids from the CURRENT BLOCK.

Selectable ids in the CURRENT BLOCK:
- n1
- n2

Pick up to <pick_limit> candidates most relevant to the query, best first.
\end{lstlisting}

\noindent\textbf{Expected tool call}
\begin{lstlisting}
rank({
  "ranked_ids": ["n1"],
  "done": true
})
\end{lstlisting}

The implementation maps local aliases such as \texttt{n1} back to real node ids
before splitting ranked ids into the top-candidate result set and the next
frontier.

\end{document}
